\documentclass[12pt,a4paper,UTF8]{ctexart}
\usepackage[a4paper,top=22mm,bottom=22mm,left=22mm,right=22mm,headheight=15pt]{geometry}
\newcommand{\Needspace}[1]{\par\begingroup\dimen0=#1\relax\dimen2=\pagegoal\advance\dimen2 by-\pagetotal\ifdim\dimen0>\dimen2\newpage\fi\endgroup}
\newcommand{\needspace}[1]{\par\begingroup\skip0=\lastskip\vskip-\skip0\dimen0=#1\relax\advance\dimen0 by\skip0\vskip\dimen0\penalty-100\vskip-\dimen0\vskip\skip0\penalty9999\endgroup}
\usepackage{amsmath,amssymb,mathtools,bm}
\usepackage{graphicx,xcolor,booktabs,longtable,array,enumitem,fancyhdr,caption,float,microtype,setspace}
\usepackage[hidelinks,unicode]{hyperref}
\usepackage[most]{tcolorbox}
\definecolor{ContestBlue}{HTML}{17365D}
\hypersetup{pdftitle={第053套},pdfauthor={}}
\setstretch{1.24}
\setlength{\parindent}{2em}
\setlength{\parskip}{0.35em}
\setlength{\emergencystretch}{2em}
\setlist{itemsep=0.35em,topsep=0.4em}
\captionsetup{font=small,labelfont=bf,justification=centering}
\setcounter{secnumdepth}{0}
\pagestyle{fancy}
\fancyhf{}
\fancyhead[L]{\small 第053套}
\fancyfoot[C]{\thepage}
\renewcommand{\headrulewidth}{0.4pt}
\ctexset{
  section={format=\Large\bfseries\color{ContestBlue},beforeskip=1.2em,afterskip=0.7em},
  subsection={format=\large\bfseries\color{ContestBlue},beforeskip=1.0em,afterskip=0.5em},
  subsubsection={format=\normalsize\bfseries,beforeskip=0.8em,afterskip=0.35em}
}
\usepackage{calc}
\newcounter{none}
\let\originalsubsection\subsection
\renewcommand{\subsection}{\needspace{6\baselineskip}\originalsubsection}
\makeatletter
\newcommand{\tightlist}{\setlength{\itemsep}{0pt}\setlength{\parskip}{0pt}}
\newcommand{\maxwidth}{\ifdim\Gin@nat@width>\linewidth\linewidth\else\Gin@nat@width\fi}
\newcommand{\maxheight}{\ifdim\Gin@nat@height>0.72\textheight 0.72\textheight\else\Gin@nat@height\fi}
\newcommand{\pandocbounded}[1]{\begingroup\setkeys{Gin}{width=\maxwidth,height=\maxheight,keepaspectratio}#1\endgroup}
\makeatother
\begin{document}
\begin{titlepage}
  \centering
  \vspace*{58mm}
  {\Huge\bfseries\color{ContestBlue} 第053套\par}
  \vfill
  \vspace*{28mm}
\end{titlepage}
\section{第 1 题}\label{ux7b2c-1-ux9898}

有两个星球------蔡子星和黄伯星，均可视为质点，质量均为
\(m\)。两个星球分别被 \(x\) 星人和 \(y\) 星人限制为只能沿 \(x\) 轴和
\(y\) 轴自由运动。初始时两星相距 \(L\)：蔡子星位于 \(x\)
轴正半轴且初速度为零，黄伯星位于原点并具有沿 \(y\) 轴负方向的初速度
\(u_0\)。\(u_0\)
足够小，使两星不会彼此分离。另有一个质量可忽略的``中点星人''，始终位于两星连线的中点。万有引力常量为
\(G\)。求：

（1）当蔡子星和黄伯星的坐标分别为 \((x,0)\)、\((0,y)\) 时，求
\(\ddot x,\ddot y\)；

（2）中点星人的加速度；

（3）中点星人的运动可等效为绕一个虚拟中心天体运动，求虚拟中心天体的质量；

（4）中点星人离原点的最近距离和最远距离；

（5）在题图坐标系中写出中点星人的轨道方程。

\begin{figure}[H]
\centering
\includegraphics[width=0.643\linewidth,height=\textheight,keepaspectratio,alt={图1-1　两星的约束运动与连线中点}]{../figures/question-01.pdf}
\caption*{图1-1　两星的约束运动与连线中点}
\end{figure}

\newpage

\section{第 2 题}\label{ux7b2c-2-ux9898}

半径为 \(R\) 的水平圆盘上放置一根半径 \(r=nR\)、长度
\(mR\)（\(m>1\)）、质量 \(M\)
均匀分布的刚性滚筒。圆盘可绕中心竖直轴无摩擦转动。滚筒中心轴为
\(AA'\)：\(A\) 端光滑铰接在圆盘轴线上，离盘面的距离为 \(r\)；\(A'\)
端由竖直光滑限位器限制，使轴线 \(AA'\)
始终位于纸面内。外力矩使滚筒绕自身轴以恒定角速度 \(\omega_0\) 转动。

（1）圆盘允许微小形变，中心不变形、边缘形变最大，形变随半径线性变化；正压力微元
\(\Delta N\) 正比于形变。滚筒与圆盘间为干性摩擦
\(f=\mu\Delta N\)。求圆盘稳定角速度之比 \(\omega_1/\omega_0\)；

（2）在上述稳定状态，求 \(A\) 点所受水平力与 \(Mg\) 的比值 \(F/(Mg)\)；

（3）若接触处单位长度所受摩擦力正比于相对速度，求圆盘稳定角速度之比
\(\omega_2/\omega_0\)。

\begin{figure}[H]
\centering
\includegraphics[width=0.823\linewidth,height=\textheight,keepaspectratio,alt={图2-1　圆盘、滚筒、铰点与限位器的空间关系}]{../figures/question-02.pdf}
\caption*{图2-1　圆盘、滚筒、铰点与限位器的空间关系}
\end{figure}

\newpage

\section{第 3 题}\label{ux7b2c-3-ux9898}

一个长圆柱形绝缘薄壳长为 \(l\)、半径为 \(R\)，满足
\(l\gg R\)，可绕水平中心轴转动，转动惯量为 \(I\)。质量为 \(m\)
的小球通过轻绳绕在圆柱上，从静止释放，绳与圆柱不打滑。薄壳均匀带总电荷
\(Q\)；转动后产生磁场和感生电场。忽略辐射与狭义相对论效应，静电力常量为
\(k\)，真空磁导率为 \(\mu_0\)。求：

（1）已知某时刻角加速度 \(\beta_1\)、角速度 \(\omega_1\)，求 \(r_1=R/2\)
处磁感应强度 \(B_1\)；

（2）求 \(r_2=2R\) 处磁感应强度 \(B_2\)；

（3）求 \(r_1=R/2\) 处电场大小 \(E_1\)；

（4）求 \(r_2=2R\) 处电场大小 \(E_2\)；

（5）求实际角加速度 \(\beta\)；

（6）小球下降高度 \(H\) 时装置的总动能 \(E_k\)；

（7）求系统磁场能量 \(E_{e1}\)；

（8）判断用磁场能量密度、等效电感储能、能流密度时间积分三种方法能否得到
\(mgH-E_k\)。

\begin{figure}[H]
\centering
\includegraphics[width=0.623\linewidth,height=\textheight,keepaspectratio,alt={图3-1　带电圆柱薄壳与下落小球}]{../figures/question-03.pdf}
\caption*{图3-1　带电圆柱薄壳与下落小球}
\end{figure}

\newpage

\section{第 4 题}\label{ux7b2c-4-ux9898}

图（a）、（b）为两种由相同元件组成的交流电桥。电感阻抗
\(\widetilde Z_L=i\omega L\)，电容阻抗
\(\widetilde Z_C=1/(i\omega C)\)。

（1）按图（a）接入 \(U(t)=U_0\cos\omega t\)，其中
\(\omega=2\omega_0=2/\sqrt{LC}\)，求稳态干路电流振幅 \(I_0\)；

（2）求干路电流相对电源的相位差 \(\varphi_0\)；

（3）按图（b）接入
\(U(t)=U_0\cos\omega_0t\)，\(\omega_0=1/\sqrt{LC}\)，求电桥中间电阻内向上电流的振幅
\(I_1\)；

（4）求该电流相对电源的相位差 \(\varphi_1\)；

（5）求图（b）的干路电流振幅 \(I_2\)；

（6）求干路电流相对电源的相位差 \(\varphi_2\)。

\begin{figure}[H]
\centering
\includegraphics[width=0.9\linewidth,height=\textheight,keepaspectratio,alt={图4-1　同臂与交叉臂交流电桥}]{../figures/question-04.pdf}
\caption*{图4-1　同臂与交叉臂交流电桥}
\end{figure}

\newpage

\section{第 5 题}\label{ux7b2c-5-ux9898}

某种粒子的速率概率密度为 \[
f(v)=Av^2\left(-a+\frac bv+\frac c{v^2}\right),
\qquad a,b,c>0,
\] 并且在允许区间内 \(f(v)>0\)。已知 \(a=3,b=2,c=1\)。求：

（1）最概然速率 \(v_p\)；

（2）截止速率 \(v_j\)，即 \(f(v_j)=0\) 且 \(v>v_j\) 后不再有粒子；

（3）归一化常数 \(A\)；

（4）平均速率 \(\bar v\) 与方均根速率
\(v_{\rm rms}=\sqrt{\overline{v^2}}\)。

\begin{figure}[H]
\centering
\includegraphics[width=0.68\linewidth,height=\textheight,keepaspectratio,alt={图5-1　有限区间内的速率概率密度}]{../figures/question-05.pdf}
\caption*{图5-1　有限区间内的速率概率密度}
\end{figure}

\newpage

\section{第 6 题}\label{ux7b2c-6-ux9898}

两个理想平面镜相对放置，光子频率为
\(\nu_0\)，在两镜间来回反射。把光子视为满足相对论能量动量关系的粒子。初始镜间距为
\(l_0\)，左镜固定，右镜以 \(u_0\ll c\)
向左匀速靠近。所有结果均在地面参考系中给出。求：

（1）光子经两次反射（每面镜各一次）后的频率 \(\nu_1\)；

（2）镜间距缩小到 \(l_0/2\) 时的频率 \(\nu_{1/2}\)；

（3）镜间距为 \(xl_0\) 时，光子对运动镜的平均作用力。

\begin{figure}[H]
\centering
\includegraphics[width=0.728\linewidth,height=\textheight,keepaspectratio,alt={图6-1　固定镜与匀速靠近的运动镜}]{../figures/question-06.pdf}
\caption*{图6-1　固定镜与匀速靠近的运动镜}
\end{figure}

\newpage

\section{第 7 题}\label{ux7b2c-7-ux9898}

水中的发光体位于距器皿壁 \(x\)
处，在器皿外壁贴有一个平凸薄透镜。透镜在空气中的焦距为
\(f\)，透镜和器皿壁均很薄；水的折射率 \(n_w=4/3\)，玻璃折射率
\(n_G=3/2\)。物体位于主光轴上，以透出光方向为正。求：

（1）像到壁距离 \(y\) 的表达式；

（2）随 \(x\) 改变，物像同侧、异侧的规律；

（3）位置关系改变的较大临界值 \(x_0\)；

（4）\(x=f\) 时的像距 \(y\)；

（5）\(x=f\) 时的放大倍数 \(\beta\)；

（6）若透镜改贴在器皿内壁，求一般物距\(x\)对应的像距\(y'\)；

（7）求内贴时的一般放大倍数\(\beta'(x)\)。

\begin{figure}[H]
\centering
\includegraphics[width=0.8\linewidth,height=\textheight,keepaspectratio,alt={图7-1　水---玻璃---空气界面上的平凸薄透镜}]{../figures/question-07.pdf}
\caption*{图7-1　水---玻璃---空气界面上的平凸薄透镜}
\end{figure}

\newpage

\section{第 8 题}\label{ux7b2c-8-ux9898}

一个可绕中心竖直轴转动的装置由两端各一匹质量为 \(m\) 的木马和一根长度为
\(L\)、质量为 \(M\) 的均匀绝缘细杆组成。杆上均匀带电，线电荷密度为
\(\lambda\)。空间存在一个关于转轴对称、半径为
\(r\)（\(r<L/2\)）、方向竖直向上的均匀圆柱形磁场。初始磁场为
\(B_0\)，随后减弱至零并保持不变；忽略阻力。

（1）若 \(B(t)=B_0-kt\)，求加速转动过程中的总力矩；

（2）求最终角速度；

（3）若 \[
B(t)=\frac{B_0}{6}f(t),
\qquad
f(t)=6-\left(\frac t\tau\right)^3
+3\left(\frac t\tau\right)^2
-5\left(\frac t\tau\right),
\] 其中 \(\tau=1\ \mathrm s\)，求磁场减到零所需时间；

（4）在第（3）问条件下，求一匹木马在加速过程中切向加速度的最小值；

（5）在第（3）问条件下，求最终角速度。

\begin{figure}[H]
\centering
\includegraphics[width=0.699\linewidth,height=\textheight,keepaspectratio,alt={图8-1　带电转杆穿过轴对称变化磁场}]{../figures/question-08.pdf}
\caption*{图8-1　带电转杆穿过轴对称变化磁场}
\end{figure}
\end{document}
